\subsubsection{Ablation on Calibration Dataset}
\label{sec:ablation_calibration}

\begin{wraptable}{r}{0.4\linewidth}
    \vspace{-1.2em}
    \centering
    \caption{Cross-dataset calibration coverage. Dataset A is WikiText-2 train.}
    \label{tab:ablation_calibration}
    \scriptsize
    \setlength{\tabcolsep}{5pt}
    \renewcommand{\arraystretch}{1.08}
    \begin{tabular}{llcc}
        \toprule
        \textbf{Eval set} & \textbf{Domain} & \textbf{A$\rightarrow$B} & \textbf{B$\rightarrow$B} \\
        \midrule
        WikiText-2 & Language & 99.65\% & 99.65\% \\
        HumanEval & Code & 99.66\% & 99.66\% \\
        GSM8K & Math & 99.64\% & 99.64\% \\
        MMLU & Knowledge & 99.63\% & 99.63\% \\
        PTB & Language & 99.55\% & 99.58\% \\
        \bottomrule
    \end{tabular}
    \vspace{-1.7em}
\end{wraptable}

SplitZip uses a calibrated top-16 exponent codebook to avoid online histogram construction.
To test whether this codebook is dataset-specific, we calibrate on WikiText-2 train and evaluate coverage on datasets from different domains.
Table~\ref{tab:ablation_calibration} compares cross-dataset calibration, denoted A$\rightarrow$B, with oracle calibration on each target dataset, denoted B$\rightarrow$B.

The calibrated exponent set generalizes well.
Coverage remains between 99.55\% and 99.66\% across all evaluated datasets, and matches the oracle result on WikiText-2 test, HumanEval~\cite{chen2021humaneval}, GSM8K~\cite{cobbe2021gsm8k}, and MMLU~\cite{hendrycks2021mmlu}.
Even on PTB~\cite{zhou2020ptb}, the gap is only 0.03 percentage points.
This shows that high-frequency KV-cache exponents are stable across domains, enabling SplitZip to use a fixed offline codebook.
